% META: {"id": "1NSI-WEB-M1", "chapitre": "1NSI-WEB-IHM", "type_objet": "methode", "capacites": ["P-IHM-02"], "capacites_codes": ["C3"], "methodes": ["M1"], "mode_creation": "ex_nihilo", "sources_inspiration": [], "status": "generated"}
\begin{fichemethode}{M1}{Analyser et modifier ce qui s'exécute au clic d'un bouton}
\textbf{Quand l'utiliser ?} Quand on te donne une page Web et qu'on te demande ce que fait
un bouton, ou de changer son comportement. Le chemin est toujours le même, et il se remonte
dans un seul sens : du bouton vers la fonction.

\textbf{Pas à pas :}
\begin{enumerate}
  \item \textbf{Trouve le composant} dans le HTML. Repère son identifiant :
    \lstinline{<button id="convertir">}. C'est cet identifiant qui relie les deux fichiers.
  \item \textbf{Trouve l'écouteur} dans le JavaScript : cherche cet identifiant. Il apparaît
    dans un \lstinline{addEventListener("click", ...)}, ou dans un attribut
    \lstinline{onclick} du HTML.
  \item \textbf{Lis la fonction associée}, et elle seule. Note trois choses : d'où elle lit
    ses données, ce qu'elle calcule, où elle écrit son résultat.
  \item \textbf{Modifie} — et une seule chose à la fois. Change le calcul, ou la source, ou
    la destination, pas les trois.
  \item \textbf{Recharge la page et reclique.} Le navigateur garde souvent l'ancienne version
    en mémoire : si rien ne change, recharge en forçant avant de conclure que ta
    modification est fausse.
\end{enumerate}

\textbf{Exemple :}

\begin{console}
<!-- page.html -->
<input id="celsius" type="text">
<button id="convertir">Convertir</button>
<p id="resultat"></p>
\end{console}

\begin{console}
// page.js
const bouton = document.getElementById("convertir");
bouton.addEventListener("click", function () {
    const c = Number(document.getElementById("celsius").value);
    const f = c * 9 / 5 + 32;
    document.getElementById("resultat").textContent = f + " F";
});
\end{console}

Lecture : la fonction lit le champ \lstinline{celsius} \emph{(source)}, calcule
$c \times 9/5 + 32$ \emph{(traitement)}, écrit dans le paragraphe \lstinline{resultat}
\emph{(destination)}. Pour convertir en Kelvin, seule la ligne du calcul change.

\textbf{Trois questions, trois endroits :}
\begin{center}
\begin{tabular}{l|l|l}
question & où regarder & ce qu'on y trouve \\
\hline
quel composant ? & HTML & \texttt{id="convertir"} \\
quel événement ? & JavaScript & \texttt{addEventListener("click", …)} \\
quel traitement ? & corps de la fonction & source, calcul, destination \\
\end{tabular}
\end{center}

\textbf{Pièges :}
\begin{itemize}
  \item \textbf{Chercher la fonction avant le composant.} Une page contient plusieurs
    écouteurs ; c'est l'identifiant qui dit lequel te concerne.
  \item Confondre \lstinline{value} et \lstinline{textContent} :
    \lstinline{value} lit ce qu'un utilisateur a saisi dans un champ,
    \lstinline{textContent} lit ou écrit le texte affiché d'un élément.
  \item \textbf{Oublier que \lstinline{value} renvoie du texte.} Sans
    \lstinline{Number(...)}, \lstinline{"20" * 9} donne bien $180$ mais
    \lstinline{"20" + 32} donne \lstinline{"2032"} : l'addition concatène.
  \item Modifier trois choses puis tester : quand le résultat est faux, on ne sait plus
    laquelle est en cause.
  \item \textbf{Mettre des parenthèses au gestionnaire.}
    \lstinline{ajouter_gestionnaire("bouton", compter())} \textbf{appelle} la fonction tout de
    suite et associe au bouton la \emph{valeur renvoyée} : le bouton devient inerte. Sans
    parenthèses, on confie la fonction elle-même.
\end{itemize}

\textbf{Vérifier :} clique avec une valeur dont tu connais la réponse — $0\,^\circ$C doit
donner $32$, $100\,^\circ$C doit donner $212$. Puis clique avec le champ vide, pour voir ce
que la page fait d'une saisie absente.

\textbf{S'entraîner :} exercices C3 du chapitre.
\end{fichemethode}
